% Part:sets-functions-relations
% Chapter: sets
% Section: reduction

\documentclass[../../../include/open-logic-section]{subfiles}

\begin{document}

\olfileid{sfr}{siz}{red}

\olsection{تقلیل}

\begin{editorial}
  ایہ حصہ تقلیل نال گنتی وار فہرست نہ بنن جوگ ہون نوں ثابت کردا اے،
  تے \olref[nen]{sec} دے نتیجیاں نال میل کھاندا اے۔ اک متبادل، ذرا
  مختصر روپ جیہڑا \olref[nen-alt]{sec} دے نتیجیاں نال میل کھاندا اے،
  \olref[red-alt]{sec} وچ دتا گیا اے۔
\end{editorial}

اسی قطری دلیل نال وکھایا سی کہ $\Pow{\PosInt}$، !!{nonenumerable} اے۔
ساڈے کول پہلے ای اک ثبوت سی کہ $\Bin^\omega$، یعنی $0$ تے $1$ دیاں
ساریاں لامتناہی لڑیاں دا سیٹ، !!{nonenumerable} اے۔ $\Pow{\PosInt}$
دے !!{nonenumerable} ہون نوں ثابت کرن دا اک ہور طریقہ ایہ اے:
وکھاؤ کہ \emph{جے $\Pow{\PosInt}$، !!{enumerable} اے تاں
  $\Bin^\omega$ وی !!{enumerable} اے}۔ کیوں جے اسی جاندے آں کہ
$\Bin^\omega$، !!{enumerable} نہیں، ایس لئی $\Pow{\PosInt}$ وی نہیں
ہو سکدا۔ اینوں اک مسئلے نوں دوجے مسئلے ول \emph{تقلیل دینا} آکھیا
جاندا اے---ایتھے اسی $\Bin^\omega$ دی گنتی وار فہرست بناؤن دے مسئلے
نوں $\Pow{\PosInt}$ دی گنتی وار فہرست بناؤن دے مسئلے ول تقلیل دیندے
آں۔ پچھلے مسئلے دا حل---$\Pow{\PosInt}$ دی اک گنتی وار فہرست---پہلے
مسئلے دا حل---$\Bin^\omega$ دی اک گنتی وار فہرست---دے دیوے گا۔

اسی سیٹ~$B$ دی گنتی وار فہرست بناؤن دے مسئلے نوں سیٹ~$A$ دی گنتی
وار فہرست بناؤن دے مسئلے ول کِویں گھٹائیے؟ اسی~$A$ دی اک گنتی وار
فہرست نوں~$B$ دی گنتی وار فہرست وچ بدلّن دا طریقہ دیندے آں۔ ایہ کرن
دا سب توں سوکھا طریقہ اک !!a{surjective} فنکشن
$f\colon A \to B$ تعریف کرنا اے۔ جے $x_1$، $x_2$، \dots{}،
~$A$ دی گنتی وار فہرست اے، تاں $f(x_1)$، $f(x_2)$، \dots{}،
~$B$ دی گنتی وار فہرست ہووے گی۔ ساڈی صورت وچ اسی ہر ہدف قدر لین
والا فنکشن $f\colon \Pow{\PosInt} \to \Bin^\omega$ لبھ رہے آں۔

\begin{prob}
وکھاؤ کہ جے کوئی !!{injective} فنکشن $g\colon B \to A$ موجود اے،
تے~$B$، !!{nonenumerable} اے، تاں~$A$ وی ایویں ای اے۔ ایہ وکھا کے
کرو کہ تسیں~$g$ نوں~$A$ دی اک گنتی وار فہرست توں~$B$ دی گنتی وار
فہرست بناؤن لئی کِویں ورت سکدے او۔
\end{prob}

\begin{proof}[{\olref[nen]{thm:nonenum-pownat}} دا تقلیلی ثبوت]
فرض کرو کہ $\Pow{\PosInt}$، !!{enumerable} ہوندا، تے ایس لئی اوہدی اک
گنتی وار فہرست $Z_{1}$، $Z_{2}$، $Z_{3}$، \dots ہوندی۔

فنکشن $f \colon \Pow{\PosInt} \to \Bin^\omega$ ایویں تعریف کرو کہ
$f(Z)$ اوہ لڑی $s$ ہووے جیہدے لئی $s(n) = 1$ عین اوس ویلے جدوں
$n \in Z$، تے ہور صورت وچ $s(n) = 0$ ہووے۔
% OLSIZ-004 ماخذ دی درستی: نتیجے والی لڑی دا ناں s اے؛ انگریزی عبارت وچ غیر متعین k والے s_k تے s_k(n) ورتے گئے سن۔
ایہ صاف طور تے اک فنکشن تعریف کردا اے، کیوں جے جدوں وی
$Z \subseteq \PosInt$ ہووے، ہر $n \in \PosInt$ یا تاں $Z$ دا
!!a{element} اے یا نہیں۔ مثال لئی، مثبت جفت عدداں دے سیٹ
$2\PosInt = \{2, 4, 6, \dots\}$ نوں لڑی $010101\dots$ ول، خالی سیٹ
نوں $0000\dots$ ول، تے خود سیٹ $\PosInt$ نوں $1111\dots$ ول بھیجیا
جاندا اے۔

ایہ !!{surjective} وی اے: $0$ تے $1$ دی ہر لڑی مثبت صحیح عدداں دے
کسے سیٹ نال میل کھاندی اے، یعنی اوس سیٹ نال جیہدے رکن اوہ صحیح عدد
نیں جیہڑے لڑی وچ~$1$ والے مقاماں دے مطابق نیں۔ ہور ٹھیک طور تے،
فرض کرو $s \in \Bin^\omega$۔ $Z \subseteq \PosInt$ نوں ایویں تعریف کرو:
\[
Z = \Setabs{n \in \PosInt}{s(n) = 1}
\]
فیر $f(Z) = s$، جیویں~$f$ دی تعریف نوں ویکھ کے جانچیا جا سکدا اے۔

ہُن فہرست اُتے غور کرو:
\[
f(Z_1), f(Z_2), f(Z_3), \dots
\]
کیوں جے $f$، !!{surjective} اے، ایس لئی $\Bin^\omega$ دا ہر رکن
~$f$ دی کسے دلیل لئی قدر لازماً اے، تے سو فہرست وچ لازماً آوے گا۔
ایس لئی ایہ فہرست~$\Bin^\omega$ دے ہر رکن دی گنتی وار
فہرست اے۔

سو جے $\Pow{\PosInt}$، !!{enumerable} ہوندا، تاں $\Bin^\omega$ وی
!!{enumerable} ہوندا۔ پر $\Bin^\omega$، !!{nonenumerable} اے
(\olref[nen]{thm:nonenum-bin-omega})۔ ایس لئی $\Pow{\PosInt}$،
!!{nonenumerable} اے۔
\end{proof}

\begin{explain}
تقلیل کس پاسے جاندی اے، ایس بارے بھلیکھا پینا سوکھا اے۔ مثال لئی،
کوئی !!a{surjective} فنکشن $g \colon \Bin^\omega \to B$ ایہ
\emph{نہیں} ثابت کردا کہ $B$، !!{nonenumerable} اے۔ (فنکشن
$g \colon \Bin^\omega \to \Bin$ اُتے غور کرو جیہڑا
$g(s) = s(1)$ نال تعریف اے، یعنی اوہ فنکشن جیہڑا $0$ تے $1$ دی
اک لڑی نوں اوہدے پہلے !!{element} ول بھیجدا اے۔ ایہ !!{surjective}
اے، کیوں جے کجھ لڑیاں $0$ توں تے کجھ $1$ توں شروع ہوندیاں نیں۔ پر
$\Bin$ متناہی اے۔) ایہ وی یاد رکھو کہ فنکشن~$f$ دا !!{surjective}
ہونا لازمی اے، نہیں تاں دلیل نہیں چلدی: $f(x_1)$، $f(x_2)$،
\dots{} وچ~$B$ دے سارے !!{element}s دے شامل ہون دی ضمانت نہیں
ہو سکدی۔ مثال لئی،
\[
h(n) = \underbrace{000\dots0}_{\text{$n$ $0$'s}}111\dots
\]
% OLSIZ-005 ماخذ دی درستی: ہر نتیجہ n صفراں توں بعد اکاں دی لامتناہی پونچھ رکھدا اے؛ انگریزی نمائش صرف متناہی صفر لکھ کے اوہنوں لامتناہی لڑی والے ہدف وچ رکھدی سی۔
اک فنکشن $h\colon \PosInt \to
\Bin^\omega$ تعریف کردا اے، پر $\PosInt$، !!{enumerable} اے۔
\end{explain}

\begin{prob}
تقلیلی دلیل نال وکھاؤ کہ مثبت صحیح عدداں دیاں جوڑیاں دے سارے
\emph{سیٹاں دے} سیٹ، !!{nonenumerable} اے۔
\end{prob}

\begin{prob}\label{sfr:siz:red:prob:nat-nat}
  تقلیلی دلیل نال وکھاؤ کہ سیٹ~$X$، جیہڑا سارے فنکشناں
  $f\colon \Nat \to \Nat$ دا سیٹ اے، !!{nonenumerable} اے (اشارا:
  $X$ توں~$\Bin^\omega$ ول
  اک ہر ہدف قدر لین والا فنکشن دیو۔)
\end{prob}

\begin{prob}
وکھاؤ کہ $\Nat^\omega$، یعنی قدرتی عدداں دیاں لامتناہی لڑیاں دا سیٹ،
تقلیلی دلیل نال !!{nonenumerable} اے۔
\end{prob}

\begin{prob}
$P$ مثبت صحیح عدداں دے سیٹ توں سیٹ $\{0\}$ ول فنکشناں دا سیٹ ہووے،
تے $Q$ مثبت صحیح عدداں دے سیٹ توں سیٹ $\{0\}$ ول
\emph{جزوی} فنکشناں دا سیٹ ہووے۔ وکھاؤ کہ~$P$، !!{enumerable} اے
تے~$Q$ نہیں۔ (اشارا: $\Bin^\omega$ دی گنتی وار فہرست بناؤن دے مسئلے
نوں~$Q$ دی گنتی وار فہرست بناؤن ول گھٹاؤ۔)
\end{prob}

\begin{prob}
$S$ مثبت صحیح عدداں دے سیٹ توں سیٹ \{0,1\} ول سارے !!{surjective}
فنکشناں دا سیٹ ہووے، یعنی~$S$ سارے !!{surjective}
~$f\colon \PosInt \to \Bin$ اُتے مشتمل اے۔ وکھاؤ کہ~$S$،
!!{nonenumerable} اے۔
\end{prob}

\begin{prob}
وکھاؤ کہ سارے حقیقی عدداں دا سیٹ~$\Real$، !!{nonenumerable} اے۔
\end{prob}

\end{document}
